\documentclass[12pt]{article}
\usepackage{amsmath,amsfonts,fullpage,amsthm,soul,xcolor}

\newtheorem*{theorem}{Theorem}
\newtheorem*{fact}{Fact}
\newtheorem*{goal}{Goal}
\newtheorem*{idea}{Idea}
\newtheorem*{defn}{Definition}
\newtheorem*{ex}{Example}

%Style section
\setlength{\textheight}{10in}
\setlength{\textwidth}{7in}
%\hoffset=-0.75in
 \pagestyle{plain}
 \setlength{\topmargin}{0in}
 \setlength{\headsep}{0in}
% \setlength{\oddsidemargin}{0in}
% \setlength{\evensidemargin}{0in}

\begin{document}


\pagestyle{empty} %use this to get rid of page numbers

\noindent
{Math 166   \hskip 2.025 truein  Names:\ \hrulefill}

\noindent
%\vskip 0.3truein
%\bigskip
%\noindent
{WS 3 - Forming groups  \hskip 1.5 truein  Section Number:\ \hrulefill\\ + area between curves}

\ 
\hrule\ \\
\noindent
\noindent
%{\bf Please start here:}
%\begin{itemize}
%\item Decide on a team name.
%\item Navigate to the Excel spreadsheet in your recitation channel in Microsoft Teams and fill out the information about your team and team members.
%\end{itemize}
%\hrule\ \\
%\\
\noindent
Find a group of 3-4 students who you are willing to work with for the next several weeks.
\begin{enumerate}
\item What is one thing you all have in common?
\vfill
\item What are your strengths as a member of this team?  Answer separately for each member.
\vfill
\vfill
\item How will you set up your group to ensure everybody contributes effectively?  (For example, this could involve establishing group roles, or deciding on communication norms, or something else.)
\vfill
\item How will you check up on your group dynamics throughout the semester?
\vfill
%\item How will you avoid mismatched expectations or contributions between group members?
%\vfill
\item How will you internally resolve problems that arise between group members?
\vfill
%\item Come up with a schedule for who will be the scribe for each week for work done in class and in the recitation section (to be turned in on Gradescope). Some of you may have to scribe more than others, so you may take this into consideration.
%\vfill
%\begin{center}
%\begin{tabular}{|c|c|c|c|}
%	\hline
%	Week starting & Scribe & Week Starting & Scribe \\ \hline
%	Week 2 & \hspace{1.5 in} & Week 7 & \hspace{1.5 in} \\ \hline
%	Week 3 & & Week 8 & \\ \hline
%	Week 4 & & Week 9 & \\ \hline
%	Week 5 & & Week 10 & \\ \hline
%	Week 6 & & & \\ \hline
%\end{tabular}
%\end{center}
%\vfill
\end{enumerate}

\newpage

\begin{goal} \rm
Review u-substitution and use the integral to find the area of various regions. In each problem, be sure to draw a graph of the functions.
\end{goal}
%\smallskip

\begin{enumerate}
%\setcounter{enumi}{-1}

\item Evaluate $\displaystyle \int \cot(2x) \; dx$.  (Hint: Re-write cotangent in terms of sine and cosine.)

\vspace{1in}

\item Evaluate $\displaystyle \int_0^4 \frac{x}{\sqrt{x^2 + 1}} \, dx$.

\vspace{1in}


\item Find the area enclosed between $y=\sin(x)$, $y=\cos(x)$, $x=0$, and $x=\pi$.
%over the interval $[0,\pi]$.

\vspace{2in}

\item Find the area of the region bounded by $y=e^x$, $y=14-e^x$, and the $y$-axis.



\vspace{2in}

\item Find the area of the region enclosed between $x=y^2-5$ and $x=3-y^2$. 

\vspace{2in}

\end{enumerate}




\end{document}
